\documentclass[12pt,a4paper]{article}

\usepackage{array}
\usepackage{booktabs}
\usepackage[margin=1in]{geometry}

\title{Creating Tables in LaTeX}
\author{Abhishek Narayan Hathwala \\ MNIT Jaipur \\ Student ID: 2025UCP1029}
\date{\today}

\begin{document}

\maketitle

\section{Introduction}

LaTeX provides the tabular environment to create
structured tables in documents.

\section{Basic Table}

\begin{center}
\begin{tabular}{|c|c|c|}
\hline
Name & Roll Number & Branch \\
\hline
Abhishek & 2025UCP1029 & CSE \\
Rahul & 2025UCP1030 & CSE \\
Aman & 2025UCP1031 & ECE \\
\hline
\end{tabular}
\end{center}

\section{Table with Different Alignments}

The following table demonstrates left, center,
and right alignment.

\begin{center}
\begin{tabular}{|l|c|r|}
\hline
Subject & Marks & Grade \\
\hline
Mathematics & 90 & A \\
Physics & 85 & A \\
Computer Science & 95 & A+ \\
\hline
\end{tabular}
\end{center}

\section{Table with Caption and Label}

\begin{table}[h]
\centering
\caption{Student Information}
\label{tab:student}

\begin{tabular}{|c|c|c|}
\hline
Name & Department & Year \\
\hline
Abhishek & CSE & 2025 \\
Rahul & ECE & 2025 \\
Aman & ME & 2025 \\
\hline
\end{tabular}

\end{table}

Table \ref{tab:student} shows sample student information.

\section{Professional Table Using Booktabs}

\begin{table}[h]
\centering
\caption{Academic Performance}

\begin{tabular}{lcc}
\toprule
Subject & Marks & Grade \\
\midrule
Mathematics & 90 & A \\
Physics & 85 & A \\
Computer Science & 95 & A+ \\
\bottomrule
\end{tabular}

\end{table}

\section{Conclusion}

This practical demonstrates how to create tables
using LaTeX, apply different alignments, add captions,
and create professional tables using the booktabs package

\end{document}